%! Graphing Exponential Functions and Transformations — Unit 7, Lesson 7.2 — Homework
\documentclass[10pt]{article}
\usepackage{algebra2-article}
\usepackage{algebra2-boxes}
\newcommand{\TallMath}[1]{$\displaystyle #1\rule[-1.4em]{0pt}{3.2em}$}
\newcommand{\lgrid}[4]{%
  \draw[step=1, linegray!40, very thin] (#1,#3) grid (#2,#4);
  \draw[->, semithick, charcoal] (#1,0)--(#2,0) node[below right, font=\scriptsize]{$x$};
  \draw[->, semithick, charcoal] (0,#3)--(0,#4) node[left, font=\scriptsize]{$y$};}
% #1=a  #2=ln(b)  #3=h  #4=k  #5=xmin  #6=xmax
% ln2 = 0.693147   ln(1/2) = -0.693147
\newcommand{\expgraph}[6]{\draw[very thick, forest, domain=#5:#6, samples=140, smooth]
  plot (\x,{(#1)*exp((#2)*((\x)-(#3)))+(#4)});}
\newcommand{\hasym}[3]{\draw[navy, dashed, semithick] (#1,#3)--(#2,#3);}

\begin{document}

\pageheader{Unit 7, Lesson 7.2}{Homework}
\namedateperiod

\vspace{0.10in}

\begin{notesbox}{Practice}
Find the asymptote before the range, every time. Give ranges in interval notation.

\begin{enumerate}[leftmargin=1.9em, itemsep=12pt]
  \item \textbf{Straight from the equation.} Fill in every cell.
    \begin{center}
    \renewcommand{\arraystretch}{1.6}
    \begin{tabularx}{\linewidth}{@{}l X X X X@{}}
    \hline
    \textbf{Function} & \textbf{Asymptote} & \textbf{Range} & \textbf{$y$-intercept} & \textbf{Growth or decay?} \\
    \hline
    $y=2^{x}+4$ & \blank{1.8cm} & \blank{2.2cm} & \blank{1.8cm} & \blank{1.8cm} \\
    $y=5(3)^{x}$ & \blank{1.8cm} & \blank{2.2cm} & \blank{1.8cm} & \blank{1.8cm} \\
    $y=\left(\frac{1}{2}\right)^{x}-2$ & \blank{1.8cm} & \blank{2.2cm} & \blank{1.8cm} & \blank{1.8cm} \\
    $y=-4^{x}$ & \blank{1.8cm} & \blank{2.2cm} & \blank{1.8cm} & \blank{1.8cm} \\
    $y=2^{\,x-3}$ & \blank{1.8cm} & \blank{2.2cm} & \blank{1.8cm} & \blank{1.8cm} \\
    \hline
    \end{tabularx}
    \end{center}
    Which rows in that table have an $x$-intercept? \blank{6.0cm}

  \item \textbf{Say what moved.} Each function below is a transformation of the parent $y=3^{x}$.
    \begin{center}
    \renewcommand{\arraystretch}{1.6}
    \begin{tabularx}{\linewidth}{@{}l X l@{}}
    \hline
    \textbf{Function} & \textbf{The move(s), in words} & \textbf{New asymptote} \\
    \hline
    $y=3^{x}-7$ & \blank{6.0cm} & \blank{1.8cm} \\
    $y=3^{\,x+4}$ & \blank{6.0cm} & \blank{1.8cm} \\
    $y=\frac{1}{2}\cdot 3^{x}$ & \blank{6.0cm} & \blank{1.8cm} \\
    $y=3^{-x}$ & \blank{6.0cm} & \blank{1.8cm} \\
    $y=-3^{x}+2$ & \blank{6.0cm} & \blank{1.8cm} \\
    \hline
    \end{tabularx}
    \end{center}

  \item \textbf{Match.} Each graph is a transformation of $y=2^{x}$; the dashed line is its asymptote.
    \begin{center}
    \begin{tikzpicture}[scale=0.27]
      % W: y = 2^x + 1
      \begin{scope}
        \lgrid{-3}{4}{-6}{8}
        \begin{scope}\clip (-3,-6) rectangle (4,8); \expgraph{1}{0.693147}{0}{1}{-3}{4} \end{scope}
        \hasym{-3}{4}{1}
        \node[charcoal, font=\small] at (0.5,-7.6) {\textbf{W}};
      \end{scope}
      % X: y = 2^x - 4
      \begin{scope}[xshift=9.5cm]
        \lgrid{-3}{4}{-6}{8}
        \begin{scope}\clip (-3,-6) rectangle (4,8); \expgraph{1}{0.693147}{0}{-4}{-3}{4} \end{scope}
        \hasym{-3}{4}{-4}
        \node[charcoal, font=\small] at (0.5,-7.6) {\textbf{X}};
      \end{scope}
      % Y: y = 2^(-x)
      \begin{scope}[xshift=19cm]
        \lgrid{-3}{4}{-6}{8}
        \begin{scope}\clip (-3,-6) rectangle (4,8); \expgraph{1}{-0.693147}{0}{0}{-3}{4} \end{scope}
        \hasym{-3}{4}{0}
        \node[charcoal, font=\small] at (0.5,-7.6) {\textbf{Y}};
      \end{scope}
      % Z: y = -2^x + 4
      \begin{scope}[xshift=28.5cm]
        \lgrid{-3}{4}{-6}{8}
        \begin{scope}\clip (-3,-6) rectangle (4,8); \expgraph{-1}{0.693147}{0}{4}{-3}{4} \end{scope}
        \hasym{-3}{4}{4}
        \node[charcoal, font=\small] at (0.5,-7.6) {\textbf{Z}};
      \end{scope}
    \end{tikzpicture}
    \end{center}
    \vspace{2pt}
    $y=2^{x}+1$ is \blank{1.0cm} \quad $y=2^{x}-4$ is \blank{1.0cm} \quad
    $y=2^{-x}$ is \blank{1.0cm} \quad $y=-2^{x}+4$ is \blank{1.0cm}

    \vspace{3pt}
    Two of those four graphs cross the $x$-axis. Name them, and give the $x$-intercept of each by
    reading the graph. \blank{7.0cm}

  \item \textbf{Backwards.} Each curve has the form $y=a\,b^{x}+k$. Use the asymptote first.
    \begin{center}
    \begin{minipage}[c]{0.30\linewidth}
    \centering
    \begin{tikzpicture}[scale=0.34]
      \lgrid{-3}{4}{-5}{9}
      \begin{scope}\clip (-3,-5) rectangle (4,9); \expgraph{2}{0.693147}{0}{-3}{-3}{4} \end{scope}
      \hasym{-3}{4}{-3}
      \fill[charcoal] (0,-1) circle (3.5pt) node[left, font=\scriptsize]{$(0,-1)$};
      \fill[charcoal] (1,1) circle (3.5pt) node[right, font=\scriptsize]{$(1,1)$};
      \node[navy, font=\scriptsize, left] at (-1.2,-3.7) {$y=-3$};
    \end{tikzpicture}\\[2pt]{\small (a)}
    \end{minipage}
    \hspace{0.02\linewidth}
    \begin{minipage}[c]{0.30\linewidth}
    \centering
    \begin{tikzpicture}[scale=0.34]
      \lgrid{-2}{5}{-1}{12}
      \begin{scope}\clip (-2,-1) rectangle (5,12); \expgraph{6}{-0.693147}{0}{1}{-2}{5} \end{scope}
      \hasym{-2}{5}{1}
      \fill[charcoal] (0,7) circle (3.5pt) node[right, font=\scriptsize]{$(0,7)$};
      \fill[charcoal] (1,4) circle (3.5pt) node[right, font=\scriptsize]{$(1,4)$};
      \node[navy, font=\scriptsize, right] at (2.6,1.7) {$y=1$};
    \end{tikzpicture}\\[2pt]{\small (b)}
    \end{minipage}
    \hspace{0.02\linewidth}
    \begin{minipage}[c]{0.32\linewidth}
    \small
    (a) $k=$ \blank{1.0cm} \; $a=$ \blank{1.0cm} \; $b=$ \blank{1.0cm} \\[4pt]
    \hspace*{1.0em} $y=$ \blank{2.6cm} \\[10pt]
    (b) $k=$ \blank{1.0cm} \; $a=$ \blank{1.0cm} \; $b=$ \blank{1.0cm} \\[4pt]
    \hspace*{1.0em} $y=$ \blank{2.6cm}
    \end{minipage}
    \end{center}

  \item \textbf{Error analysis.} Name the error, then give the correct answer.

    \textbf{(a)} Asked to describe $y=5^{\,x-2}$, Marcus wrote: ``The asymptote is $y=-2$ and the
    range is $(-2,\infty)$.'' \writelines{2}

    \textbf{(b)} Tara says $y=-6^{x}$ and $y=6^{-x}$ are both decay functions ``because of the
    negative.'' Fix both halves of that claim. \writelines{2}
\end{enumerate}
\end{notesbox}

\begin{extensionbox}
\textbf{A soda out of the refrigerator.} A can of soda at $38^{\circ}$F is set on a counter in a
room held at $72^{\circ}$F. Its temperature after $t$ minutes is
\[
  T(t) = 72 - 34(0.9)^{t}.
\]
\begin{center}
\renewcommand{\arraystretch}{1.25}
\begin{tabular}{|c|c|c|c|c|c|c|}
\hline
$t$ (min) & $0$ & $5$ & $10$ & $15$ & $20$ & $30$ \\ \hline
$T(t)$ ($^{\circ}$F) & $38$ & $51.9$ & $60.1$ & $65.0$ & $67.9$ & $70.6$ \\ \hline
\end{tabular}
\end{center}
\begin{enumerate}[label=(\alph*), leftmargin=1.8em, itemsep=5pt]
  \item Identify $a$, $b$, and $k$. Why is $a$ \emph{negative} here, when the soda is warming up?
  \item What is the horizontal asymptote, and what does it mean about this soda?
  \item Give the range of $T$ as an algebraic function, then give the temperatures the soda actually
        takes for $t\ge0$. Explain why the two differ, and why one endpoint is a bracket.
  \item Because $a<0$ and $k>0$ have opposite signs, this function \emph{does} have an
        $x$-intercept. Solving $34(0.9)^{t}=72$ would require $(0.9)^{t}=2.118\ldots$ --- explain why
        that forces $t$ to be \emph{negative}, and say why the $x$-intercept is therefore irrelevant
        to this story.
  \item At what time is the soda exactly $72^{\circ}$F? Justify with the asymptote.
\end{enumerate}
\writelines{6}
\end{extensionbox}

\begin{spiralbox}
\textbf{Coming next --- Lesson 7.3: Solving Exponential Equations with a Common Base.} In problem 3
you found the $x$-intercept of $y=2^{x}-4$ by \emph{looking} at the graph. Next lesson you solve for
it: $2^{x}-4=0$ becomes $2^{x}=4$, and since $4$ is already a power of $2$, the exponents can simply
be matched. That one-to-one move runs entirely on Unit 6's exponent machinery --- $8^{x}=4$,
$3^{x}=\frac{1}{27}$, $\sqrt[3]{2}=2^{1/3}$ --- and every answer can be checked as an intersection on
a graph like the ones above. Notice, though, what it \emph{cannot} do: the coffee equation
$132(0.85)^{t}=32$ from the Activity and the soda's $(0.9)^{t}=2.118$ above have no common base at
all. Hold on to those two.
\end{spiralbox}

\end{document}
